\section{欧洲监管与主权：把课堂时点和当前规则分开}

学生问到 EU AI Act 时，Lample 的回答反映了当时状态：高层透明度要求已经可见，但具体需要披露什么、技术模板如何执行仍在讨论。这个判断不能原封不动写成 2026 年现状。欧盟委员会随后明确，general-purpose AI（GPAI）模型提供者的透明度与版权相关义务自 2025 年 8 月 2 日起适用；在该日期前已投放市场的模型有至 2027 年 8 月 2 日的过渡期，达到系统性风险门槛的模型还有额外安全义务。

\begin{warningbox}{监管时间注记，不构成法律意见}
本节只解释课程中的基础设施影响。具体模型是否属于 GPAI、谁是 provider、何时算 placed on the market，以及开源例外或系统性风险义务如何适用，需要结合欧盟委员会现行指引和专业法律意见。规则和解释会继续更新。
\end{warningbox}

\subsection{为什么 privacy 与 reliability 会改变市场边界}

金融、医疗、国防和保险客户不只关心模型分数，还关心数据是否离开边界、服务依赖是否可控、版本是否能冻结、故障时能否回滚。即使外部 API 的工程质量更高，组织仍可能因为责任边界和业务连续性选择 private deployment。读图时应把 privacy、reliability 与 control 当作三个独立轴：组织可能只缺其中一个，也可能愿意为三者同时获得更强保证而承担更高运营成本。

\lecturefigure{15-privacy-reliability.png}{敏感组织购买的是 privacy、reliability 与 control 的组合。}{本地字幕 35:40--37:10；概念重绘。}

\begin{knowledgebox}{读图：主权不是“所有东西都自己造”}
主权更接近可替换性和决策权：能否选择模型，迁移数据，审计依赖，冻结版本，在供应商退出时保持服务。组织可以使用外部模型或开源资产，同时通过标准接口、数据导出、模型评测和多供应商策略降低锁定。
\end{knowledgebox}

\teachervoice{课堂中的欧洲动机同时包含人才、产业和主权。Lample 也明确表示自己并非监管专家，并把许多判断限定为“too early to say”。讲义保留这种限定，避免把创业者观点包装成政策事实。}

\subsection{本章小结}

EU AI Act 必须按日期理解：课堂时点的技术细节仍不确定，之后 GPAI obligations 已进入适用阶段。对系统设计而言，监管和主权最终会落到数据 lineage、文档、评测、安全控制和可审计部署。

